Terug in haar heden,
los van haar verleden,

van de incarnatie
Maria naar Emmy,

weg uit de zodiak,
waar de angst aanbrak,

naast de jager ontwaakt,
stond ze weer poedelnaakt.

Bloot was in de oertijd
wellicht een gangbaar feit,

maar in een beschaving
vaak een overtreding.

Het was privébezit.
De jager ontving dit

niet als gezichtsverlies
en had geen zichtverlies.

Kalm werd ze bekeken.
Ze hoorde hem spreken

over zijn vulva vangst.
Dat bekeek hij het langst.

De drankgeur uit zijn mond,
de ijskoude moerasgrond,

maar bovenal naaktheid
trotseerden haar doodsstrijd

tegen uitlevering.
Zoals Dogville los ging

door Grace haar vluchtpaniek,
nam Emmy haar tragiek.

Haar brein, veel te duchten,
wilde daarvoor vluchten.

Ze broedde een brief uit,
die over de schavuit,

wat er aan vooraf ging
en haar motivering

voor haar overgave,
aan haar zuster Emma

helder zou voorleggen.
Wat ze niet kon zeggen

aan de woeste jager,
haar hete belager.

Vaak met was verzegeld
werd een brief geregeld

door een koerier gebracht.
Niet dat zij had verwacht

dat haar mentaal schrijven
op schrift zou beschrijven,

aan haar oudere zus,
wat gebeurde aldus.

De denkbrief aan Emma
had volgens Derrida

niet het onzegbare
in het openbare

domein tot taal gebracht,
maar wel een spoor getracht

te trekken van haar wond
met een gesloten mond.

Al konden de meesten
van de arme geesten

niet lezen of schrijven,
door de kloosterwijven,

die Emmy had ontmoet,
leerde zij dat juist goed

en zou het, afgezien
van haar huidig afzien,

ook zeker verstuurd zijn,
in pagina’s vol pijn.

Haar woorden verklaarden,
als feiten aanvaarden

de mensen hun strekking,
die globaal als volgt ging.

In de brief aan Emma
beschreef Emmy hoe ma

kutkamde op een keer
tegen hun ouweheer

balend als een stekker.
“Waarom ben je gekker

op je paard dan op mij”,
was wat ze bijtend zei,

staand in een paardenvijg.
“Als ik mijn paard bestijg

dan wil hij graag rijden!”.
Emmy kon aansnijden

het verstandshuwelijk,
dat ze als gruwelijk

als kind had ervaren,
met deze bezwaren.

“Iemand die ijdel zeurt,
vraagt om een flinke beurt.”

Moeder moest onderop
voor een hop paardje hop.

Moeder zocht caritas,
maar kreeg cupiditas,

een lagere liefde
wat hun binding kliefde.

“Zoals ma pa niet wou,
onderga ik de kou,

die de jager uitstraalt.
Werd ik maar aangehaald

door een knappe jager.
De heer, schoon en mager,

die me eerder snapte
en gretig toehapte.

Werd ik maar bekeken
door de jonge leken,

die in hun onkunde,
hun aftrekken gunde,

als stroopsmerend gebaar,
zonder enig gevaar.”

Dit was, haar hoop ten spijt,
niet haar realiteit.

Haar pa kwam tot inkeer.
De verrukte landheer

vernam van zijn voordeel.
Men bracht naar zijn kasteel

zowel Emmy’s mama
als haar zuster Emma.

“Emma je eerste brief
was begripvol en lief.

Je begreep mijn onrust.
De vermeende wellust,

die je had ontboden,
beschreef je als noden

in onverbloemde stijl.
Je kwam daar voor de bijl,

omdat in deze tijd
men kijkt naar een dienstmeid

als erotisch object,
vanwege het aspect

van de autoriteit
versus noeste arbeid.

De landheer schoot als schurk
recht omhoog als een kurk,

toen de kans zich aanbood.
Dus hij jullie ontbood.

Machiavellisme
is zijn mechanisme

om zijn macht te houden,
maar jullie vertrouwden

hem desondanks toch wel.
Voor hem is het een spel.

Zijn persoonlijk gewin
geeft oppervlakkig zin

aan zijn vlakke leven.
Nemen zonder geven.

Hij verkreeg positie
door manipulatie

en de titel hertog
door sluwheid en bedrog.

Maar deze goede raad
komt helaas veel te laat.”

“Zoals ma had voorspeld
werd een meisje geveld

door een mensenjager.
Niemand werd de vrager

van wie de klos zou zijn?
Het lijkt mijn verhaallijn

kennelijk te schrijven.
Je kunt me aanwrijven

dat ik mijn lijf aanbied,
maar welk mens doet dat niet

als in een wanhoopsdaad
alles op het spel staat.

Wat mis ik het klooster
en zijn strakke rooster!

Had ik daar maar geschuild,
daar vrijheid ingeruild.”

Zo sloeg de profetie
van mama op Emmy.

Eveneens roodharig
en net meerderjarig

was de murwe balling
in moeders voorspelling.

Net als bij Jesaja
accepteerde Emma

moeder haar heilsprofeet,
die de voorspelling deed.

Emmy meende echter,
als alerte rechter,

dat de komst van Jezus
een staging was, versus

een echte profetie.
Doch twijfelde Emmy

of deze achterdocht,
gezien haar rampentocht,
wel van toepassing was,
zo staand bij het moeras;

gevlucht, meisje, gember
en op elf november.

Zo kreeg ze vulva-angst!
Verviel zij tot vulvangst?
